\section{Schwarz maps on direct sums of matrix algebras}

\begin{definition}[Kraus maps between finite sums of matrix algebras]
    \label{def:kraus_map_between_direct_sums}
    \lean{Matrix.directSumMapExtension, Matrix.IsKrausDirectSumMap}
    \leanok
    Let
    \[
        \mathcal A=\bigoplus_{i\in I}M_{n_i}(\mathbb C),
        \qquad
        \mathcal B=\bigoplus_{j\in J}M_{m_j}(\mathbb C),
    \]
    and write $\iota_{\mathcal A},\pi_{\mathcal A}$ and
    $\iota_{\mathcal B},\pi_{\mathcal B}$ for the corresponding
    block-diagonal embeddings and diagonal-block compressions.  The canonical
    full-matrix extension of $T:\mathcal A\to\mathcal B$ is
    \[
        \widehat T=\iota_{\mathcal B}\circ T\circ\pi_{\mathcal A}.
    \]
    The map $T$ is a direct-sum Kraus map when $\widehat T$ has a Kraus
    representation.
\end{definition}

% This is the finite-sum complete-positivity implication used for the
% transported composites in CPSV16, Appendix C.4, lines 1974--1995.  It does
% not establish Kraus representations for those transported maps; that is the
% tensor-specific preparation and partial-trace step tracked in issue #4438.
\begin{lemma}[Composition and trace-adjoint Schwarz for direct-sum Kraus maps]
    \label{lem:direct_sum_kraus_composition_schwarz}
    \lean{Matrix.directSumMapExtension_comp,
        Matrix.IsKrausDirectSumMap.comp,
        Matrix.isSchwarzDirectSumMap_of_directSumExtension_isSchwarzMap,
        Matrix.IsKrausDirectSumMap.directSumTraceAdjointMap_isSchwarz}
    \leanok
    \uses{def:kraus_map_between_direct_sums}
    For direct-sum Kraus maps
    $T:\mathcal A\to\mathcal B$ and $S:\mathcal B\to\mathcal C$, one has
    \[
        \widehat{S\circ T}=\widehat S\circ\widehat T,
    \]
    and $S\circ T$ is again a direct-sum Kraus map.  If an endomorphism
    $T:\mathcal A\to\mathcal A$ has an extension satisfying
    \[
        \widehat T(Y^*Y)-\widehat T(Y^*)\widehat T(Y)\succeq0,
    \]
    then $T$ satisfies the corresponding inequality in every summand.
    Finally, if $F:\mathcal A\to\mathcal A$ is a trace-preserving direct-sum
    Kraus map, then its trace adjoint satisfies, for every $X\in\mathcal A$,
    \[
        F^*(X^*X)-F^*(X^*)F^*(X)\succeq0
    \]
    in every summand.
\end{lemma}

\begin{proof}\leanok
    \uses{thm:kadison_schwarz_2positive}
    Since $\pi_{\mathcal B}\iota_{\mathcal B}=\id_{\mathcal B}$,
    \[
        \widehat S\,\widehat T
        =\iota_{\mathcal C}S\pi_{\mathcal B}
          \iota_{\mathcal B}T\pi_{\mathcal A}
        =\iota_{\mathcal C}ST\pi_{\mathcal A}
        =\widehat{S\circ T}.
    \]
    If $\widehat T(X)=\sum_i K_iXK_i^*$ and
    $\widehat S(Y)=\sum_j L_jYL_j^*$, then
    \[
        \widehat{S\circ T}(X)
          =\sum_{i,j}(L_jK_i)X(L_jK_i)^*,
    \]
    which gives the asserted Kraus representation.  For
    $Y=\iota_{\mathcal A}(X)$, the $i$-th diagonal block of the extension
    Schwarz inequality is
    \[
        \left[\widehat T(Y^*Y)-\widehat T(Y^*)\widehat T(Y)\right]_i
          =(T(X^*X))_i-(T(X^*))_i(T(X))_i\succeq0,
    \]
    which proves the descent assertion.  If $F$ preserves the total trace,
    then $\widehat F$ preserves the ordinary trace, so $(\widehat F)^*$ is
    unital.  Complete positivity implies two-positivity for
    $(\widehat F)^*$, and Kadison--Schwarz gives
    \[
        (\widehat F)^*(Y^*Y)
          -(\widehat F)^*(Y^*)(\widehat F)^*(Y)\succeq0.
    \]
    For $Y=\iota_{\mathcal A}(X)$, use
    $(\widehat F)^*=\widehat{F^*}$ and project this inequality onto the
    $i$-th diagonal summand.  The resulting block identity is
    \[
        \left[\widehat{F^*}(Y^*Y)
          -\widehat{F^*}(Y^*)\widehat{F^*}(Y)\right]_i
        =(F^*(X^*X))_i-(F^*(X^*))_i(F^*(X))_i\succeq0.
    \]
\end{proof}

\begin{theorem}[Mutually inverse CPTP maps give a star equivalence]%
    \label{thm:direct_sum_inverse_cptp_star_equiv}
    \lean{IsKrausCP.kadison_schwarz_of_map_one_eq_one}
    \lean{Matrix.IsSchwarzBetweenDirectSums}
    \lean{Matrix.IsKrausDirectSumMap.isSchwarzBetweenDirectSums}
    \lean{Matrix.IsKrausDirectSumMap.map_conjTranspose_between}
    \lean{Matrix.schwarz_equality_of_mutual_inverse_kraus_direct_sum_maps}
    \lean{Matrix.map_mul_of_mutual_inverse_kraus_direct_sum_maps}
    \lean{Matrix.starAlgEquivOfMutualInverseKrausDirectSumMaps}
    \lean{Matrix.starAlgEquivOfMutualInverseKrausDirectSumMaps_apply}
    \lean{Matrix.starAlgEquivOfMutualInverseKrausDirectSumMaps_symm_apply}
    \lean{Matrix.directSumTraceAdjointMapBetween_comp_eq_id_of_comp_eq_id}
    \lean{Matrix.directSumTraceAdjointMapBetween_map_mul_of_mutual_inverse}
    \lean{Matrix.directSumTraceAdjointStarAlgEquiv}
    \lean{Matrix.directSumTraceAdjointStarAlgEquiv_apply}
    \lean{Matrix.directSumTraceAdjointStarAlgEquiv_symm_apply}
    \leanok
    \uses{def:kraus_map_between_direct_sums,
        lem:direct_sum_rectangular_trace_adjoint}
    Let $T:\mathcal A\to\mathcal B$ and $S:\mathcal B\to\mathcal A$ be
    mutually inverse completely positive trace-preserving maps between finite
    direct sums of full matrix algebras.  Their trace adjoints are mutually
    inverse unital completely positive maps.  Moreover, for all
    $X,Y\in\mathcal B$,
    \[
        T^*(XY)=T^*(X)T^*(Y),
        \qquad T^*(X^*)=T^*(X)^*,
    \]
    and likewise for $S^*$.  Thus $T^*$ and $S^*$ determine mutually inverse
    star-algebra equivalences.  This theorem does not yet identify the simple
    summands or their dimensions.
\end{theorem}

\begin{proof}\leanok
    Complete positivity and trace preservation make both trace adjoints
    unital Schwarz maps.  For $F=T^*$ and $G=S^*$, let
    \[
        \Delta_F(X)=F(X^*X)-F(X)^*F(X)\succeq0.
    \]
    Positivity of $G$ gives $G(\Delta_F(X))\succeq0$.  The Schwarz inequality
    for $G$, together with $GF=\id$, gives
    $-G(\Delta_F(X))\succeq0$.  Hence $G(\Delta_F(X))=0$, and injectivity of
    $G$ implies $\Delta_F(X)=0$.  Polarization of this equality gives
    multiplicativity.  Positivity gives preservation of the adjoint, so the
    mutually inverse maps are star-algebra equivalences.  This supplies an
    algebraic route to the block-classification step for which
    \cite[Appendix~C.4, line~1997]{Cirac2016MPDO_arXiv} cites
    \cite[Theorem~8]{WolfPerezGarcia2010Inverse}.  The latter instead uses
    extreme states and the positive inverse to identify the blocks, then the
    Schwarz condition to exclude transposition.
\end{proof}

\begin{theorem}[Simple-summand matching for inverse CPTP maps]%
    \label{thm:direct_sum_inverse_cptp_block_matching}
    \lean{Matrix.exists_blockEquiv_dim_eq_of_mutual_inverse_kraus_direct_sum_maps}
    \leanok
    \uses{thm:direct_sum_inverse_cptp_star_equiv,
        thm:dim_preserved}
    Let $\{D_i\}_{i\in I}$ and $\{E_j\}_{j\in J}$ be positive integers, and
    set
    \[
        \mathcal A=\prod_{i\in I}M_{D_i}(\C),
        \qquad
        \mathcal B=\prod_{j\in J}M_{E_j}(\C).
    \]
    If $T:\mathcal A\to\mathcal B$ and
    $S:\mathcal B\to\mathcal A$ are mutually inverse completely positive
    trace-preserving maps, let $\Phi:\mathcal A\simeq\mathcal B$ be the
    star-algebra isomorphism obtained from their trace adjoints.  Write $I_i$
    and $J_j$ for the block ideals of $\mathcal A$ and $\mathcal B$,
    respectively.  Then there is an equivalence $\sigma:I\simeq J$ such that
    \[
        \Phi(I_i)=J_{\sigma(i)},
        \qquad
        D_i=E_{\sigma(i)}
    \]
    for every $i\in I$. This is the simple-summand matching conclusion used
    in \cite[Appendix~C.4, line~1997]{Cirac2016MPDO_arXiv}; the unitary action
    within the paired summands is a separate conclusion.
\end{theorem}

\begin{proof}
    \leanok
    \uses{thm:direct_sum_inverse_cptp_star_equiv,
        thm:dim_preserved}
    By Theorem~\ref{thm:direct_sum_inverse_cptp_star_equiv}, the rectangular
    trace adjoints determine a star-algebra isomorphism between
    $\mathcal A$ and $\mathcal B$. Theorem~\ref{thm:dim_preserved} applied to
    this isomorphism gives the equivalence~$\sigma$, proves that it carries
    $I_i$ onto $J_{\sigma(i)}$, and gives the equalities
    $D_i=E_{\sigma(i)}$.
\end{proof}

\begin{theorem}[Unitary Skolem--Noether theorem for complex matrices]%
    \label{thm:unitary_skolem_noether_matrix}
    \lean{MPSTensor.exists_unitary_conj_of_starAlgEquiv_matrix}
    \leanok
    Every star-algebra automorphism $\Phi$ of a nonzero full complex matrix
    algebra is implemented by a unitary: there is a unitary matrix $U$ such
    that
    \[
        \Phi(X)=UXU^*
    \]
    for every matrix $X$.
\end{theorem}

\begin{proof}\leanok
    \uses{thm:skolem_noether}
    By Skolem--Noether, there is an invertible matrix $P$ such that
    $\Phi(X)=PXP^{-1}$.  Preservation of the adjoint implies that $P^*P$
    commutes with every matrix, hence $P^*P=c\Id$ for a scalar $c$.  Positivity
    and invertibility give $c>0$.  Therefore $U=c^{-1/2}P$ is unitary and
    implements the same automorphism.
\end{proof}

\begin{theorem}[Unitary action within matched simple summands]%
    \label{thm:direct_sum_star_equiv_block_unitary}
    \lean{MPSTensor.blockComponentStarAlgEquiv}
    \lean{MPSTensor.matrixReindexStarAlgEquiv}
    \lean{MPSTensor.exists_unitary_block_implementers_of_starAlgEquiv_pi_matrix}
    \lean{Matrix.exists_blockEquiv_dim_eq_unitary_of_starAlgEquiv_pi_matrix}
    \lean{Matrix.exists_blockEquiv_dim_eq_unitary_of_mutual_inverse_kraus_direct_sum_maps}
    \leanok
    Retain an equivalence $\sigma:I\simeq J$ which matches the block ideals
    of a star-algebra isomorphism
    \[
        \Phi:\prod_{i\in I}M_{D_i}(\C)
        \simeq \prod_{j\in J}M_{E_j}(\C),
    \]
    together with equalities $D_i=E_{\sigma(i)}$.  Then there are unitaries
    $U_i\in M_{E_{\sigma(i)}}(\C)$ such that
    \[
        \Phi(0,\ldots,0,X,0,\ldots,0)_{\sigma(i)}
        =U_i\,\iota_i(X)\,U_i^*
    \]
    for every $X\in M_{D_i}(\C)$, where $\iota_i$ is the reindexing induced
    by the retained dimension equality.  In particular, this conclusion
    applies to the trace-adjoint star-algebra isomorphism of any mutually
    inverse completely positive trace-preserving pair.  This is the
    blockwise-unitary statement in
    \cite[Appendix~C.4, line~1997]{Cirac2016MPDO_arXiv}; it does not assert
    the later MPDO multiplicity or coefficient relations.
\end{theorem}

\begin{proof}\leanok
    \uses{thm:unitary_skolem_noether_matrix,
        thm:direct_sum_inverse_cptp_star_equiv,
        thm:dim_preserved}
    The restriction of $\Phi$ to a matched block is a star-algebra
    isomorphism.  After reindexing by $D_i=E_{\sigma(i)}$, it is an
    automorphism of a full complex matrix algebra.  Skolem--Noether writes
    this automorphism as conjugation by an invertible matrix $P$.  Preservation
    of the adjoint implies that $P^*P$ commutes with every matrix and is
    therefore a positive scalar multiple of the identity.  Dividing $P$ by
    the positive square root of this scalar produces the required unitary.
\end{proof}

\begin{lemma}[Total trace under unitary action on matched summands]
    \label{lem:direct_sum_block_unitary_trace_preserving}
    \lean{Matrix.isTracePreservingBetweenDirectSums_of_unitary_block_implementers}
    \leanok
    Let $\Phi:\mathcal A\simeq\mathcal B$ match the simple summands by an
    equivalence $\sigma:I\simeq J$, with $D_i=E_{\sigma(i)}$.  Suppose that
    there are unitaries $U_i$ satisfying
    \[
      \Phi(0,\ldots,0,X,0,\ldots,0)_{\sigma(i)}
      =U_i\,\iota_i(X)\,U_i^*.
    \]
    Then $\Phi$ preserves the total block trace.
\end{lemma}

\begin{proof}\leanok
    The image of a matrix supported on the $i$-th summand is supported on the
    $\sigma(i)$-th summand.  Unitary conjugation and reindexing preserve the
    trace, so
    \[
      \sum_j\tr
        \bigl(\Phi(0,\ldots,0,X,0,\ldots,0)_j\bigr)
      =\tr(X).
    \]
    Summing this identity over the source summands proves the result.
\end{proof}

\begin{theorem}[Unitary action of an invertible CPTP map]
    \label{thm:direct_sum_inverse_cptp_forward_unitary}
    \lean{Matrix.%
      exists_blockEquiv_dim_eq_unitary_forward_of_mutual_inverse_kraus_direct_sum_maps}
    \leanok
    Let $T:\mathcal A\to\mathcal B$ and $S:\mathcal B\to\mathcal A$ be
    mutually inverse completely positive trace-preserving maps between finite
    products of nonzero full matrix algebras.  There are an equivalence
    $\sigma:I\simeq J$, equalities $D_i=E_{\sigma(i)}$, and unitaries
    $U_i\in M_{E_{\sigma(i)}}(\C)$ such that
    \[
      T(0,\ldots,0,X,0,\ldots,0)
      =(0,\ldots,0,U_i\,\iota_i(X)\,U_i^*,0,\ldots,0)
    \]
    with the nonzero entry on the right in position $\sigma(i)$, and
    \[
      S(0,\ldots,0,Y,0,\ldots,0)
      =(0,\ldots,0,\iota_i^{-1}(U_i^*YU_i),0,\ldots,0)
    \]
    with the entries on the left and right in positions $\sigma(i)$ and $i$,
    respectively.  These identities hold for every $i\in I$,
    $X\in M_{D_i}(\C)$, and $Y\in M_{E_{\sigma(i)}}(\C)$.
\end{theorem}

\begin{proof}\leanok
    \uses{thm:direct_sum_inverse_cptp_star_equiv,
      thm:direct_sum_star_equiv_block_unitary,
      lem:direct_sum_block_unitary_trace_preserving,
      lem:direct_sum_trace_adjoint_star_equiv_inverse,
      lem:direct_sum_rectangular_trace_adjoint}
    Let $\Phi=S^*:\mathcal A\simeq\mathcal B$ be the star-algebra
    isomorphism obtained from the inverse pair.  Theorem~
    \ref{thm:direct_sum_star_equiv_block_unitary} supplies $\sigma$, the
    dimension equalities, and the unitaries for $\Phi$.  Lemma~
    \ref{lem:direct_sum_block_unitary_trace_preserving} shows that $\Phi$
    preserves the total trace, hence Lemma~
    \ref{lem:direct_sum_trace_adjoint_star_equiv_inverse} gives
    $\Phi^*=\Phi^{-1}$.  On the other hand, involutivity of the trace adjoint
    gives $\Phi^*=S$.  Therefore $S=\Phi^{-1}$, and the two inverse laws imply
    $T=\Phi$.  The unitary formula for $\Phi$ gives the asserted formula for
    $T$, while the inverse formula for $S=\Phi^{-1}$ uses the same summand
    matching and the same unitaries.
\end{proof}

% Scope restriction documented in
% docs/paper-gaps/wolf_theorem6_14_fixed_point_projection_gap.tex.
\begin{theorem}[Weighted block form of the adjoint Ces\`aro projection on full support]%
    \label{thm:adjoint_mean_ergodic_projection_weighted_block_form}
    \lean{IsPositiveMap.exists_block_densities_of_adjoint_meanErgodicProjection}
    \leanok
    \uses{thm:positive_unital_adjoint_mean_ergodic_retraction,
        thm:fixedPointsStarSubalgebra,
        thm:positive_retraction_finite_star_algebra}
    Let $T:\MN{D}\to\MN{D}$ be positive and trace preserving, suppose that
    $T^*$ satisfies the Schwarz inequality, and suppose that there is a
    positive definite matrix $\rho$ satisfying $T(\rho)=\rho$.  Let $P_T$ be
    the mean-ergodic projection of $T$.  There are positive integers $d_k,m_k$,
    a unitary $U$, and density matrices $\sigma_k\in\MN{m_k}$ such that
    \[
        U^*\Fix(T^*)U
        =\bigoplus_k\Id_{m_k}\otimes\MN{d_k}
    \]
    and, for every $A\in\MN{D}$,
    \[
        P_T^*(A)=U\left(\bigoplus_k\Id_{m_k}\otimes
        \tr_{m_k}\!\left(
          (\sigma_k\otimes\Id_{d_k})(U^*AU)_{kk}\right)\right)U^*.
    \]
    This is the full-support conditional-expectation step in the proof of
    \cite[Theorem~6.14]{Wolf2012Quantum}.
\end{theorem}

\begin{proof}
    \leanok
    \uses{thm:positive_unital_adjoint_mean_ergodic_retraction,
        thm:fixedPointsStarSubalgebra,
        thm:positive_retraction_finite_star_algebra}
    The map $P_T^*$ is a positive unital idempotent map whose range is
    $\Fix(T^*)$.  For $A\in\Fix(T^*)$, the
    Schwarz inequality gives
    $T^*(A^*A)\geq T^*(A)^*T^*(A)=A^*A$, while
    \[
        \tr\!\left(\rho T^*(A^*A)\right)
        =\tr\!\left(T(\rho)A^*A\right)
        =\tr(\rho A^*A).
    \]
    Since $\rho$ is positive definite, these relations imply
    $T^*(A^*A)=A^*A$.  Hence $\Fix(T^*)$ is a $*$-subalgebra,
    and $P_T^*$ is a positive retraction onto it.  The blockwise
    classification of positive retractions gives the unitary, the matrix
    dimensions, the density matrices $\sigma_k$, and the displayed formula.
\end{proof}

% Scope restriction and factor-order convention documented in
% docs/paper-gaps/wolf_theorem6_14_fixed_point_projection_gap.tex.
\begin{theorem}[Density-block form of the Ces\`aro projection on full support]%
    \label{thm:mean_ergodic_projection_density_block_form}
    \lean{IsPositiveMap.exists_block_densities_of_meanErgodicProjection}
    \leanok
    \uses{thm:adjoint_mean_ergodic_projection_weighted_block_form}
    Under the hypotheses of
    Theorem~\ref{thm:adjoint_mean_ergodic_projection_weighted_block_form},
    the mean-ergodic projection of $T$ satisfies
    \[
        P_T(B)=U\left(\bigoplus_k \sigma_k\otimes
          \tr_{m_k}((U^*BU)_{kk})\right)U^*.
    \]
    Consequently,
    \[
        \Fix(T)
        =U\left(\bigoplus_k \sigma_k\otimes\MN{d_k}\right)U^*.
    \]
    Here each summand is written on
    $\mathbb C^{m_k}\otimes\mathbb C^{d_k}$, so
    $\tr_{m_k}$ traces the first factor.  This is the
    full-support part of Equation~(6.63) in
    \cite[Theorem~6.14]{Wolf2012Quantum}; the density matrices are not yet
    asserted to be positive definite.
\end{theorem}

\begin{proof}
    \leanok
    \uses{thm:adjoint_mean_ergodic_projection_weighted_block_form}
    For one summand, cyclicity of the trace and the defining property of the
    partial trace give
    \[
      \tr\!\left(
        (\sigma_k\otimes\tr_{m_k}B)A\right)
      =\tr\!\left(
        B\bigl(\Id_{m_k}\otimes
          \tr_{m_k}((\sigma_k\otimes\Id_{d_k})A)\bigr)\right).
    \]
    Summing over the diagonal summands and applying the unitary change of
    basis proves the displayed formula by nondegeneracy of the trace pairing.
    Its range consists of the stated density blocks because
    $\tr_{m_k}(\sigma_k\otimes X)=X$.
\end{proof}

% Scope restriction documented in
% docs/paper-gaps/wolf_theorem6_14_fixed_point_projection_gap.tex.
\begin{theorem}[Density-block form of fixed points on a finite direct sum]%
    \label{thm:direct_sum_fixed_points_density_block_form}
    \lean{Matrix.IsPositiveDirectSumMap.exists_block_densities_of_fixedPoints}
    \leanok
    \uses{thm:direct_sum_full_matrix_extension_compatibility,
        thm:mean_ergodic_projection_density_block_form}
    Let
    \[
        \mathcal A=\bigoplus_{i\in I}M_{n_i}(\mathbb C)
    \]
    be a finite direct sum.  Suppose that $T:\mathcal A\to\mathcal A$ is
    positive and preserves the total trace, that its trace adjoint satisfies
    the Schwarz inequality, and that $T$ fixes a family $\rho=(\rho_i)_i$
    with every $\rho_i$ positive definite.  Then there are positive integers
    $d_k,m_k$, density matrices $\sigma_k\in M_{m_k}(\mathbb C)$, a unitary
    $U$ on $\bigoplus_i\mathbb C^{n_i}$, and a reindexing of this space by
    $\bigoplus_k(\mathbb C^{m_k}\otimes\mathbb C^{d_k})$ such that
    \[
        T(A)=A
        \quad\Longleftrightarrow\quad
        U^*\iota(A)U=\bigoplus_k \sigma_k\otimes X_k
    \]
    for some $X_k\in M_{d_k}(\mathbb C)$.  Here $\iota$ is the
    block-diagonal embedding.  This is the finite-direct-sum fixed-point step
    used in \cite[Appendix~C.4, lines~1980--1995]{Cirac2016MPDO_arXiv}; it
    does not assert the channel hypotheses for the particular sector maps in
    that argument.  This is the full-support restriction: the support
    reduction and complementary zero summand of the general theorem remain
    open.
\end{theorem}

\begin{proof}
    \leanok
    \uses{thm:direct_sum_full_matrix_extension_compatibility,
        thm:mean_ergodic_projection_density_block_form}
    Extend $T$ by diagonal compression and block-diagonal embedding.  The
    canonical extension is positive and trace preserving, its trace adjoint
    satisfies the Schwarz inequality, and the embedded family $\iota(\rho)$
    is a positive-definite fixed point.  After choosing a numbered basis of
    the finite-dimensional ambient space, apply the full-support
    density-block theorem.  Simultaneous reindexing preserves positivity,
    trace preservation, the Schwarz inequality, the trace adjoint, and
    positive definiteness.  Reindexing the resulting unitary and block
    decomposition back to $\mathcal H$ gives the displayed equivalence.
\end{proof}

Without a full-rank hypothesis on the fixed point, the corner-restricted fixed-point
$*$-algebra of Theorem~\ref{thm:cornerFixedPointsStarSubalgebra} still carries the block
structure through the compression onto the support sector; this is the
$\sum_k d_k m_k \le D$ support-sector form of the block representation in
Equation~(1.39) of~\cite{Wolf2012Quantum} invoked
by~\cite[Theorem~6.14]{Wolf2012Quantum}. A corner-restricted fixed point extended by
zero on the complement of the support sector need not be a fixed point of the ambient
map, so the corner block form does not by itself give the ambient fixed-point space a
zero-block representation.

\begin{theorem}[Block form of the corner-restricted fixed points]%
    \label{thm:cornerFixedPoints_block_form}
    \lean{Kraus.cornerFixedPoints_blockDiagonal_iff}
    \leanok
    \uses{def:support_proj, def:kraus_map_channels}
    Let $E(X) = \sum_i K_i X K_i^{\dagger}$ be a trace-preserving Kraus map on $\MN{D}$,
    with Heisenberg-picture adjoint $E^*(Y) = \sum_i K_i^{\dagger} Y K_i$, let
    $\rho \succeq 0$ satisfy
    $E(\rho) = \rho$, and let $Q$ be the support projection of $\rho$. There are a sector
    dimension $r \le D$, a number $n$ of blocks, positive dimensions $d_0, \ldots,
    d_{n-1}$ and multiplicities $m_0, \ldots, m_{n-1}$ with $\sum_k d_k m_k = r$,
    realized by an explicit identification of the index sets, and an isometry
    $W : \C^{r} \to \C^{D}$ with $W^{\dagger} W = \Id_{r}$ and $W W^{\dagger} = Q$, such
    that a matrix $Y \in \MN{D}$ satisfies $Q\,Y\,Q = Y$ and $Q\,E^*(Y)\,Q = Y$ exactly
    when
    \[
        Y \;=\; W \Bigl(\bigoplus_{k=0}^{n-1} \Id_{m_k} \otimes B_k\Bigr) W^{\dagger}
    \]
    for some matrices $B_k \in \MN{d_k}$; that is, up to reordering the two tensor
    factors of each block,
    \[
        \{Y \in Q\,\MN{D}\,Q \mid Q\,E^*(Y)\,Q = Y\} \;=\;
        W \Bigl(\bigoplus_{k=0}^{n-1} \Id_{m_k} \otimes \MN{d_k}\Bigr) W^{\dagger}.
    \]
    The right-hand side is the support-sector block representation in
    Equation~(1.39) of~\cite{Wolf2012Quantum}; the equivalence characterizes the
    corner-restricted set, not the ambient fixed-point space.
\end{theorem}

\begin{proof}
    \leanok
    \uses{thm:support_proj_fixed, thm:supported_corner_compression_isometry,
        thm:compression_on_support_posDef, thm:adjointFixedPoints_block_form}
    Since $\rho$ is a fixed point of $E$, the support projection satisfies
    $(\Id - Q) K_i Q = 0$ (Theorem~\ref{thm:support_proj_fixed}), so the compressed
    family $Q K_i Q$ is supported on the corner and trace-preserving there. Compress it
    to the support sector along the isometry $V : \C^{r} \to \C^{D}$ of
    Theorem~\ref{thm:supported_corner_compression_isometry}, which satisfies
    $V^{\dagger} V = \Id_{r}$ and $V V^{\dagger} = Q$: the compressed
    family $C$ is trace-preserving on $\C^{r}$, and the compression
    $\sigma = V^{\dagger} \rho V$ of $\rho$ is a positive definite fixed point of the
    compressed Schr\"odinger map (Theorem~\ref{thm:compression_on_support_posDef}). The
    block form of the Heisenberg-picture fixed points of $C$
    (Theorem~\ref{thm:adjointFixedPoints_block_form}) yields a unitary
    $U \in \MN{r}$, dimensions $d_k$ and multiplicities $m_k$ with
    $\sum_k d_k m_k = r$, and the equivalence between the fixed-point equation on the
    sector and the block-diagonal conjugation identity. Writing
    $\Phi(X) := V X V^{\dagger}$ for the compression isomorphism onto the corner and
    $C^*$ for the Heisenberg-picture adjoint of the compressed family, the
    corner-restricted fixed-point condition corresponds to the fixed-point equation of
    the compressed map,
    \[
        Q\,E^*\!\bigl(\Phi(X)\bigr)\,Q \;=\; \Phi(X)
        \quad\Longleftrightarrow\quad
        C^*(X) \;=\; X,
    \]
    so, with $W := V U$, which satisfies
    $W^{\dagger} W = \Id_r$ and $W W^{\dagger} = Q$, a corner element $Y$ is a
    corner-restricted fixed point exactly when $Y = W (\bigoplus_k \Id_{m_k} \otimes
    B_k) W^{\dagger}$ for some blocks $B_k$. Finally $r \le D$ because the isometry $V$
    embeds $\C^{r}$ into $\C^{D}$.
\end{proof}

\begin{theorem}[Weighted fixed points form a $*$-subalgebra]%
    \label{thm:weightedFixedPointsStarSubalgebra}
    \lean{Kraus.weightedFixedPointsStarSubalgebra}
    \leanok
    \uses{def:tp_kraus, thm:fixedPointsStarSubalgebra}
    Let $T(X) = \sum_i K_i X K_i^\dagger$ be trace-preserving, and let
    $\rho > 0$ satisfy $T(\rho) = \rho$. If
    \[
        F_T := \{X \in \MN{D} \mid T(X) = X\},
    \]
    then $\rho^{-1/2} F_T \rho^{-1/2}$ is a $*$-subalgebra of $\MN{D}$.
\end{theorem}
\begin{proof}
    \leanok
    \uses{thm:fixedPointsStarSubalgebra}
    Set $B_i = \rho^{-1/2} K_i \rho^{1/2}$. The equation $T(\rho) = \rho$
    implies that the family $\{B_i\}$ is unital, and trace preservation of $T$
    implies that $\rho$ is a positive-definite fixed point of the adjoint map
    of the gauged family. Apply Theorem~\ref{thm:fixedPointsStarSubalgebra}
    to $B$. The identity
    \[
        \sum_i B_i X B_i^\dagger = X
        \iff
        T(\rho^{1/2} X \rho^{1/2}) = \rho^{1/2} X \rho^{1/2}
    \]
    identifies the fixed points of the gauged map with
    $\rho^{-1/2} F_T \rho^{-1/2}$.
\end{proof}

For a singular fixed point the conjugation is taken with the inverse square root on the
support of $\rho$, and the conjugated set lives in the corner algebra of the support
projection. The corollary in~\cite{Wolf2012Quantum} takes a maximum-rank fixed point, for
which every fixed point of $T$ is supported on the support of $\rho$; the following two
statements take an arbitrary positive semidefinite fixed point and restrict to the fixed
points supported on the support of $\rho$. Theorem~\ref{thm:maximalSupport_fixedPoint}
below supplies a fixed point whose support carries every fixed point, and at such a fixed
point the restriction disappears
(Theorem~\ref{thm:maximalSupport_weightedCorner}).
% Maintainer note: the restriction history is recorded in
% docs/paper-gaps/wolf_cor67_maximal_support_restriction.tex.

\begin{theorem}[Weighted corner fixed points form a $*$-subalgebra]%
    \label{thm:weightedCornerFixedPointsStarSubalgebra}
    \lean{Kraus.weightedCornerFixedPointsStarSubalgebra}
    \leanok
    \uses{def:tp_kraus, def:support_proj, rem:corner_algebra_star_instances}
    Let $T(X) = \sum_i K_i X K_i^\dagger$ be trace-preserving, let $\rho \succeq 0$
    satisfy $T(\rho) = \rho$, and let $Q$ be the support projection of $\rho$. The corner
    elements $Y \in Q\,\MN{D}\,Q$ such that $\sqrt{\rho}\, Y \sqrt{\rho}$ is a fixed
    point of $T$ form a $*$-subalgebra of the corner algebra $Q\,\MN{D}\,Q$.
\end{theorem}

\begin{proof}
    \leanok
    \uses{thm:weightedFixedPointsStarSubalgebra,
        thm:supported_corner_compression_isometry, thm:compression_on_support_posDef,
        thm:support_proj_fixed}
    Closure under addition, scalars, and conjugation is linearity together with the
    stability of fixed points under conjugate transposition; the corner unit $Q$ belongs
    to the set because $\sqrt{\rho}\, Q \sqrt{\rho} = \rho$ is a fixed point, the support
    projection absorbing the square root: $Q \sqrt{\rho} = \sqrt{\rho} = \sqrt{\rho}\,Q$,
    since $(\Id - Q)\sqrt{\rho}\,\bigl((\Id - Q)\sqrt{\rho}\bigr)^{\dagger} =
    (\Id - Q)\rho(\Id - Q) = 0$.

    For closure under multiplication, compress to the support sector: since $\rho$ is a
    fixed point of $T$, the support projection satisfies $(\Id - Q) K_i Q = 0$
    (Theorem~\ref{thm:support_proj_fixed}), the compressed family $Q K_i Q$ is
    trace-preserving on the corner, and the compression along the isometry
    $V : \C^{r} \to \C^{D}$ of Theorem~\ref{thm:supported_corner_compression_isometry},
    with $V^{\dagger} V = \Id_r$ and $V V^{\dagger} = Q$, produces a trace-preserving
    family $C$ on $\C^{r}$ whose Schr\"odinger map has the positive definite fixed point
    $\sigma = V^{\dagger} \rho V$ (Theorem~\ref{thm:compression_on_support_posDef}). The
    square root compresses along the isometry,
    \[
        \sqrt{\sigma} \;=\; V^{\dagger} \sqrt{\rho}\, V,
    \]
    by uniqueness of the positive square root, since
    $(V^{\dagger}\sqrt{\rho}V)(V^{\dagger}\sqrt{\rho}V)
    = V^{\dagger}\sqrt{\rho}\,Q\sqrt{\rho}V = \sigma$. Hence for corner-supported $Y$,
    \[
        \sqrt{\sigma}\,(V^{\dagger} Y V)\,\sqrt{\sigma}
        \;=\; V^{\dagger}\,(\sqrt{\rho}\, Y \sqrt{\rho})\,V,
    \]
    and the fixed-point equations correspond: for corner-supported $W$ the compressed
    map evaluates to
    \[
        C\bigl(V^{\dagger} W V\bigr) \;=\; V^{\dagger}\,Q\,T(W)\,Q\,V
        \;=\; V^{\dagger}\,T(W)\,V,
    \]
    since $T(W)$ is again corner-supported, so $T(W) = W$ exactly when the compressed
    map fixes $V^{\dagger} W V$. The set therefore corresponds, element by element, to the weighted
    fixed points of the compressed family with respect to the positive definite fixed
    point $\sigma$, which form a $*$-subalgebra by
    Theorem~\ref{thm:weightedFixedPointsStarSubalgebra}. Closure under multiplication
    transports back through the correspondence: a corner-supported $Y_j$ satisfies
    $Y_j Q = Y_j$ (right-multiply $Q Y_j Q = Y_j$ by $Q$ and use $Q^2 = Q$), so
    \[
        (V^{\dagger} Y_1 V)(V^{\dagger} Y_2 V)
        \;=\; V^{\dagger}\, Y_1 (V V^{\dagger})\, Y_2\, V
        \;=\; V^{\dagger}\, (Y_1 Q)\, Y_2\, V
        \;=\; V^{\dagger}\,(Y_1 Y_2)\,V,
    \]
    and the product of two members corresponds to the product of their compressed
    images, which lies in the compressed $*$-subalgebra.
\end{proof}

Combined with the conjugation correspondence of
Theorem~\ref{thm:weightedCorner_sqrt_surjective}, the carrier of this $*$-subalgebra
realizes
\[
    \rho^{-1/2}\,\{X \in \MN{D} \mid T(X) = X,\ Q X Q = X\}\,\rho^{-1/2},
\]
with the inverse square root taken on the support of $\rho$: the singular case of
Corollary~6.7 of~\cite{Wolf2012Quantum}, restricted to the fixed points supported on
the support of $\rho$.

\begin{theorem}[Conjugation by the square root is onto the corner-supported fixed
    points]%
    \label{thm:weightedCorner_sqrt_surjective}
    \lean{Kraus.exists_weightedCorner_sqrt_eq_of_fixedPoint}
    \leanok
    \uses{def:tp_kraus, def:support_proj}
    In the setting of Theorem~\ref{thm:weightedCornerFixedPointsStarSubalgebra}, every
    fixed point $X$ of $T$ with $Q X Q = X$ arises as $X = \sqrt{\rho}\, Y \sqrt{\rho}$
    for a corner-supported $Y$ with $\sqrt{\rho}\, Y \sqrt{\rho}$ fixed by $T$.
\end{theorem}

\begin{proof}
    \leanok
    \uses{thm:supported_corner_compression_isometry, thm:compression_on_support_posDef,
        thm:support_proj_fixed}
    Compress to the support sector as in the proof of
    Theorem~\ref{thm:weightedCornerFixedPointsStarSubalgebra} and set
    \[
        Y \;=\; V\,\bigl(\sqrt{\sigma}^{-1}\,(V^{\dagger} X V)\,
            \sqrt{\sigma}^{-1}\bigr)\,V^{\dagger},
    \]
    where $\sigma = V^{\dagger} \rho V$ is positive definite, so $\sqrt{\sigma}$ is
    invertible. Then $Y$ is corner-supported, and, writing
    $Z = \sqrt{\sigma}^{-1}(V^{\dagger} X V)\sqrt{\sigma}^{-1}$ for the compressed
    middle factor, the transport identity displayed above gives
    \[
        V^{\dagger}\bigl(\sqrt{\rho}\,Y\sqrt{\rho}\bigr)V
        \;=\; \sqrt{\sigma}\,\bigl(V^{\dagger} Y V\bigr)\,\sqrt{\sigma}
        \;=\; \sqrt{\sigma}\,Z\,\sqrt{\sigma}
        \;=\; V^{\dagger} X V,
    \]
    hence
    $\sqrt{\rho}\, Y \sqrt{\rho} = Q X Q = X$, both sides being corner-supported. In
    particular $\sqrt{\rho}\, Y \sqrt{\rho}$ is fixed by $T$.
\end{proof}

The support hypothesis on the fixed points is discharged by exhibiting a fixed point
whose support carries every fixed point. The first ingredient is that Loewner domination
transfers the support.

\begin{lemma}[The transfer map of a trace-preserving family is a channel]%
    \label{lem:isChannel_transferMap}
    \lean{Kraus.isChannel_transferMap}
    \leanok
    \uses{def:channel, def:tp_kraus}
    Let $K_0, \ldots, K_{d-1}$ be a trace-preserving Kraus family. The map
    $E(X) = \sum_i K_i X K_i^{\dagger}$ is a quantum channel.
\end{lemma}
\begin{proof}\leanok
    \uses{thm:transfer_channel}
    The trace-preservation hypothesis is the normalization
    $\sum_i K_i^{\dagger} K_i = \mathbf 1$ of the channel form of the transfer map
    (Theorem~\ref{thm:transfer_channel}).
\end{proof}

\begin{lemma}[Scalar Loewner domination transfers the support]%
    \label{lem:stationaryProj_absorb_of_le}
    \lean{Kraus.stationaryProj_absorb_of_le_smul}
    \leanok
    \uses{def:support_proj}
    Let $\rho, P \succeq 0$ and $c\in\C$ with $P \preceq c\rho$, and let $Q$ be the
    support projection of $\rho$. Then
    \[
        Q\,P = P, \qquad P\,Q = P.
    \]
\end{lemma}

\begin{proof}
    \leanok
    \uses{def:support_proj}
    The support projection absorbs $\rho$, so $(\Id - Q)\rho = 0$ and
    \[
        0 \;\preceq\; (\Id - Q)\,P\,(\Id - Q)
        \;\preceq\; (\Id - Q)\,(c\rho)\,(\Id - Q) \;=\; 0,
    \]
    the middle inequality because conjugation by $\Id - Q$ preserves the Loewner order.
    Hence
    \[
        \bigl((\Id - Q)\sqrt{P}\bigr)\bigl((\Id - Q)\sqrt{P}\bigr)^{\dagger}
        \;=\; (\Id - Q)\,P\,(\Id - Q) \;=\; 0,
    \]
    so $(\Id - Q)\sqrt{P} = 0$, and
    \[
        (\Id - Q)\,P \;=\; \bigl((\Id - Q)\sqrt{P}\bigr)\sqrt{P} \;=\; 0,
    \]
    that is $Q P = P$; taking conjugate transposes gives $P Q = P$ since $P$ and $Q$ are
    Hermitian.
\end{proof}

\begin{theorem}[A fixed point of maximal support from bounded orbits]%
    \label{thm:maximalSupport_fixedPoint_of_bounded_orbits}
    \lean{IsPositiveMap.exists_maximalSupport_fixedPoint_of_hasBoundedOrbits}
    \leanok
    \uses{def:positive_map, def:has_bounded_orbits,
        def:mean_ergodic_projection, def:support_proj}
    Let $T:\MN{D}\to\MN{D}$ be positive and have bounded orbits, and let
    $T_\infty$ be its mean-ergodic projection.  Then
    $\rho_0=T_\infty(\Id)$ is positive semidefinite and fixed by $T$.  If
    $Q_0$ is the support projection of $\rho_0$, then every fixed point $X$
    of $T$ satisfies
    \[
        Q_0\,X\,Q_0=X.
    \]
    This is the bounded-orbit form of the maximal-support argument underlying
    \cite[Proposition~6.9]{Wolf2012Quantum}.
\end{theorem}

\begin{proof}
    \leanok
    \uses{thm:positive_mean_ergodic_projection,
        thm:mean_ergodic_bounded_orbits,
        thm:positive_map_conjTranspose, thm:psd_le_trace_smul_one,
        lem:stationaryProj_absorb_of_le}
    Positivity of the mean-ergodic projection gives $\rho_0\succeq 0$, while
    its range and idempotence give $T(\rho_0)=\rho_0$.  If $A\succeq0$, then
    $A\preceq\tr(A)\Id$; applying the positive map $T_\infty$ gives
    \[
        \tr(A)\rho_0-T_\infty(A)
        =T_\infty\!\left(\tr(A)\Id-A\right)\succeq0,
        \qquad
        T_\infty(A)\preceq\tr(A)\rho_0.
    \]
    Scalar domination therefore gives
    \[
        Q_0\,T_\infty(A)\,Q_0=T_\infty(A).
    \]
    For a fixed Hermitian matrix $H$, write $H=H^+-H^-$.  Since
    $T_\infty(H)=H$, linearity gives
    \[
        H=T_\infty(H^+)-T_\infty(H^-),
    \]
    so $Q_0HQ_0=H$.  Finally decompose an arbitrary fixed point into two
    Hermitian fixed matrices, using preservation of conjugate transpose, and
    conclude that $Q_0XQ_0=X$.
\end{proof}

\begin{theorem}[A fixed point of maximal support]%
    \label{thm:maximalSupport_fixedPoint}
    \lean{IsPositiveMap.exists_maximalSupport_fixedPoint}
    \leanok
    \uses{def:positive_map, def:trace_preserving, def:mean_ergodic_projection,
        def:support_proj}
    Let $T:\MN{D}\to\MN{D}$ be positive and trace-preserving, and let $T_\infty$ be its
    mean-ergodic projection. Then $\rho_0=T_\infty(\Id)$ is positive semidefinite and
    fixed by $T$. With $Q_0$ the support projection of $\rho_0$, every fixed point $X$
    of $T$ satisfies
    \[
        Q_0\, X\, Q_0 = X.
    \]
    This is the maximal-support property of fixed points in Section~6.4
    of~\cite{Wolf2012Quantum}.
\end{theorem}

\begin{proof}
    \leanok
    \uses{thm:positive_trace_preserving_mean_ergodic_projection,
        thm:maximalSupport_fixedPoint_of_bounded_orbits}
    Theorem~\ref{thm:positive_trace_preserving_mean_ergodic_projection} gives
    bounded orbits.  Apply
    Theorem~\ref{thm:maximalSupport_fixedPoint_of_bounded_orbits}.
\end{proof}

\begin{theorem}[Restriction to the support of a stationary positive matrix]%
    \label{thm:stationarySupport_restriction}
    \lean{IsPositiveMap.map_supported_on_fixedPoint_support}
    \leanok
    \uses{def:positive_map, def:support_proj,
        lem:stationaryProj_absorb_of_le}
    Let $T:\MN{D}\to\MN{D}$ be positive, and let
    $\rho\succeq0$ satisfy $T(\rho)=\rho$. If $Q$ is the support projection of
    $\rho$, then for every $X\in\MN{D}$,
    \[
        QXQ=X \quad\Longrightarrow\quad QT(X)Q=T(X).
    \]
\end{theorem}

\begin{proof}
    \leanok
    \uses{lem:stationaryProj_absorb_of_le}
    First suppose $A\succeq0$ and $QAQ=A$.  On the support of $\rho$, the
    matrix $\rho$ is positive definite, so $Q\preceq c\rho$ for some positive
    scalar $c$.  Moreover,
    \[
        A\preceq\tr(A)Q,
    \]
    because this follows by compressing
    $\tr(A)\Id-A\succeq0$ with $Q$.  Positivity and $T(\rho)=\rho$ therefore give
    \[
        0\preceq T(A)\preceq c\tr(A)\rho.
    \]
    Scalar domination transfers the support, hence $QT(A)Q=T(A)$.
    Write $H_1=X+X^*$ and $H_2=\mathrm{i}(X-X^*)$, so that both matrices are
    Hermitian and
    \[
        X=\tfrac{1}{2}H_1-\tfrac{\mathrm{i}}{2}H_2.
    \]
    Setting $A_j^\pm=QH_j^\pm Q$ gives positive matrices with
    $QA_j^\pm Q=A_j^\pm$ and
    \[
        X=\tfrac{1}{2}(A_1^+-A_1^-)-\tfrac{\mathrm{i}}{2}(A_2^+-A_2^-).
    \]
    Applying the positive-input conclusion to these four matrices and using
    linearity of $T$ gives the assertion.
\end{proof}

\begin{theorem}[Restriction to full-rank fixed points]%
    \label{thm:stationarySupport_compression}
    \lean{IsPositiveMap.exists_maximalSupportCompression}
    \leanok
    \uses{thm:maximalSupport_fixedPoint, thm:stationarySupport_restriction}
    Let $T:\MN{D}\to\MN{D}$ be positive and trace preserving, set
    $\rho_0=T_\infty(\Id)$, and let $Q_0$ be the support projection of
    $\rho_0$. There are a finite-dimensional Hilbert space $\mathcal H$ and an
    isometry $V:\mathcal H\to\mathbb C^D$ with
    $V^*V=\Id_{\mathcal H}$ and $VV^*=Q_0$ such that
    \[
        \widetilde T(Y)=V^*T(VYV^*)V
    \]
    is positive and trace preserving and has a positive definite fixed point.
    Moreover,
    \[
        T(VYV^*)=V\widetilde T(Y)V^*,
    \]
    and
    \[
        T(X)=X
        \quad\Longleftrightarrow\quad
        \exists\,Y:\widetilde T(Y)=Y,\quad X=VYV^*.
    \]
    These are Equations~(6.52) and~(6.51), respectively,
    of~\cite{Wolf2012Quantum}. Thus the complementary fixed-point summand
    vanishes.
\end{theorem}

\begin{proof}
    \leanok
    \uses{thm:maximalSupport_fixedPoint, thm:stationarySupport_restriction}
    Theorem~\ref{thm:maximalSupport_fixedPoint} gives $\rho_0$ and shows that
    $Q_0XQ_0=X$ for every fixed point $X$ of $T$. Choose an isometry onto the
    range of $Q_0$.
    Positivity follows from positivity of conjugation and of $T$.  The matrix
    $VYV^*$ is supported on $Q_0$, so the preceding theorem shows that its
    image is also supported there. Inserting $Q_0=VV^*$ gives
    \[
        T(VYV^*)
        =Q_0T(VYV^*)Q_0
        =VV^*T(VYV^*)VV^*
        =V\widetilde T(Y)V^*.
    \]
    By cyclicity of the trace, the isometry identity, and trace preservation of
    $T$,
    \[
        \tr(\widetilde T(Y))
        =\tr(V\widetilde T(Y)V^*)
        =\tr(T(VYV^*))
        =\tr(VYV^*)
        =\tr(Y).
    \]
    Thus $\widetilde T$ is trace preserving. The compressed point
    $\sigma_0=V^*\rho_0V$ is positive definite. Since
    $V\sigma_0V^*=Q_0\rho_0Q_0=\rho_0$,
    \[
        \widetilde T(\sigma_0)
        =V^*T(\rho_0)V
        =V^*\rho_0V
        =\sigma_0.
    \]
    If $T(X)=X$, then $Q_0XQ_0=X$, and hence
    $X=V(V^*XV)V^*$. Moreover,
    \[
        \widetilde T(V^*XV)
        =V^*T\bigl(V(V^*XV)V^*\bigr)V
        =V^*T(X)V
        =V^*XV.
    \]
    Conversely, if $\widetilde T(Y)=Y$, then the intertwining identity gives
    \[
        T(VYV^*)=V\widetilde T(Y)V^*=VYV^*.
    \]
\end{proof}

The next three theorems are generic finite-matrix facts.  They do not depend on
positive maps, fixed points, or Theorem~6.14 of~\cite{Wolf2012Quantum}.
